\documentclass[10pt]{article}

\usepackage[margin=1in]{geometry}
\usepackage{amsmath,amssymb}
\usepackage{booktabs}
\usepackage{graphicx}
\usepackage{hyperref}
\usepackage{siunitx}
\usepackage[numbers]{natbib}

% Guideline convention: flag design choices that lack a solid reason.
\newcommand{\unjustified}{\textbf{[!]}~\emph{Unjustified choice:}}

\title{Policy-Ready Data Collection on a Dual Five-Degree-of-Freedom
Rig: Timestamp-Aligned Multi-View Recording and a
Teleoperation-to-Autonomy Contract}

\author{so101\_garment project\thanks{Living document: this paper is the
canonical description of the real-world data-collection and training
platform in the \texttt{so101\_garment} repository and MUST be updated,
in the same branch, whenever the collection platform, the recording
system, the recorded stream set, or the teleoperation-to-autonomy
contract change. The repository's writing guideline governs its style.
The teleoperation system it builds on and the simulation training
environment that precedes it each have their own living paper in the
same repository.}}

\date{Last updated: \today}

\begin{document}
\maketitle

\begin{abstract}
A teleoperated demonstration only teaches autonomy when the recorded
observation and the recorded intent belong to the same instant, and when
the intent is the command the robot actually executed. Multi-view rigs
make both easy to lose: independent camera streams on shared buses arrive
at different times, proprioception arrives faster still, and the operator's
raw input is not the command the arm followed once limits and the
reachable envelope have constrained it. We describe a real-rig collection
system that timestamps every stream as it is read, aligns all streams to a
single reference time for each recorded frame, and logs the residual drift
so temporal quality is measured rather than assumed. The system records
the projected, constrained target --- the quantity a policy must reproduce
--- in both joint and end-effector form, so a joint-space or an
end-effector policy trains from the same episodes without recollection. We
state the contract that makes this data usable for autonomy: at inference
the policy's output passes back through the identical command path that
produced the labels. The platform is measured and its defining constants
travel with the dataset for reproducibility, a readiness check confirms
the rig is at a fixed collection state, and a difficulties section records
every obstacle met and how it was resolved. An experiments section grows
with this living document.
\end{abstract}

\section{Introduction}

This paper describes how we collect real-world demonstrations on a dual
five-degree-of-freedom manipulation rig so that the data can train a
policy which runs autonomously on the same platform. It first presents the
collection platform, then the recording system, then the contract that
links a teleoperated demonstration to autonomous control, then the
measurements needed to reproduce the setup, and finally an honest account
of the difficulties met. Each part forward-links to the next.

Behaviour cloning learns a map from an observation at one moment to the
intended command at that same moment~\citep{ross2011dagger}. A
demonstration is therefore only as good as two properties: the observation and the intent must belong to the
same instant, and the recorded intent must be the command the robot
actually executed. Break either property and the policy learns a jittery
or non-causal map that fails when it drives the arm on its own.

Our rig makes both properties hard to hold. Several cameras stream over
shared buses and arrive at slightly different times, while the arm's joint
signals arrive several times faster. If the recorder simply takes each
stream's most recent value when it writes a frame, the streams in that
frame are silently skewed against one another. Separately, the operator's
raw input is not what the arm followed: a rate limiter, the joint limits, a
deliberately capped gripper opening, and the reachable-workspace projection
all sit between the operator and the motors.

The first idea addresses timing. Every stream is timestamped at the instant
it is read from its device, on a monotonic clock. For each recorded frame
the collector fixes a single reference time and samples every stream at
that time, and it records the residual gap between the reference time and
each stream's nearest real sample. Temporal alignment thus becomes a
measured quantity, shown live during collection, rather than an assumption.

The second idea addresses intent. We record the projected and constrained
target --- the command the robot actually followed --- and we record it in
both joint and end-effector form. We then require that, at inference, a
policy's output passes back through the same command path that produced the
labels. A joint-space policy replays through the same limiter and gripper
cap; an end-effector policy replays through the same inverse kinematics and
envelope projection.

Our contributions are: a timestamp-aligned multi-view recorder with a
logged per-stream drift budget (Section~\ref{sec:collection}); a
teleoperation-to-autonomy contract that records the executed, constrained
target for both a joint-space and an end-effector policy
(Section~\ref{sec:contract}); a measured, reproducible platform whose
defining constants travel with the dataset (Sections~\ref{sec:platform}
and~\ref{sec:replicability}); and a candid difficulties record
(Section~\ref{sec:difficulties}).

\section{Collection platform}
\label{sec:platform}

The platform is a dual five-degree-of-freedom rig with two teleoperation
interfaces and a multi-view sensor suite. This section walks the arms, the
operator interfaces, the sensing, and the host and buses, and it ends with
the fixed state the rig must reach before collection begins.

\paragraph{Arms.} Two follower arms, each with five actuated body joints
and a gripper, sit in a fixed left and right layout. That layout is shared
with the digital twin used for earlier simulation work, so the physical
placement and the simulated placement come from the same description. Each
arm carries a camera at the wrist that looks along the gripper.

\paragraph{Operator interfaces.} The rig accepts two kinds of
teleoperation, and both produce demonstrations with identical recorded
semantics. A head-mounted controller drives the arms in Cartesian space
through an inverse-kinematics solver. A pair of leader arms, kinematically
identical to the followers, drives them directly in joint space. The choice
of interface changes how the operator moves, not what the recorder stores.

\paragraph{Sensing.} Each gripper carries two visual-tactile fingertip
cameras in addition to the wrist camera, and a colour camera overlooks the
workspace from above; a depth-capable camera may take that overhead position
instead, adding an aligned-depth image. The colour images, together with
that optional overhead depth, are the exterior views a policy sees; the
fingertip views add local contact information. All views are recorded
through one pipeline so that no stream is privileged over another.

\paragraph{Host and buses.} Every device reaches a single laptop. The
cameras and arms share two powered universal-serial-bus hubs, one per side,
and the head-mounted controller uses a direct link. Shared-bus bandwidth is
a genuine constraint when several compressed image streams run at once. We
do not assume the bus copes; instead the recorder measures each stream's
rate and dropped frames, so a saturated bus is visible rather than silent
(Section~\ref{sec:collection}).

\paragraph{Final-state readiness.} Collection is only meaningful when the
rig starts from a fixed, known state: every arm calibrated and bound to a
stable port alias, the wrist joints limited so the fingertip cables cannot
be twisted, the arms at their defined home pose, and every camera opened at
its configured resolution. Each colour camera is likewise bound to a stable
per-socket alias rather than an enumeration index, the overhead colour
camera included (a depth-capable overhead camera is pinned by its own serial
instead), so replugging a hub can never silently swap one view for another
between collection and deployment; the same bound identities drive both the
live monitor and the recorder, and the monitor shows every stream before a
session begins, including the overhead view, with its aligned depth rendered
alongside when a depth camera is fitted. The same monitor also
replays a recorded episode from the stored frames, so a session can be
reviewed exactly as it was seen during collection, and that review can be
exported as a video. A
preflight check confirms these conditions and
the host-tuning items that reduce timing jitter, and reports a single
ready-or-not verdict. The platform's fixed dimensions and camera
calibration are measured and stored so the rig can be rebuilt
(Section~\ref{sec:replicability}).

\section{Data-collection system}
\label{sec:collection}

The system produces two artefacts per session: a training-ready dataset and
a full-rate diagnostic log. This section walks the two-rate design, the
stamp-on-read rule, the reference-time alignment, the drift budget, the
recorded stream set, and the phase flag and gating.

\paragraph{Two-rate design.} A dataset carries a single frame
rate~\citep{lerobot}, so the training-ready features are recorded at one
fixed rate suited to the policy. Diagnosis, however, needs the control signals at their native, much
higher rate. We therefore keep a separate full-rate side log with
wall-clock stamps that holds every control signal, stored beside the
dataset and never fed to the policy. The dataset trains; the side log
explains.

\paragraph{Stamp on read.} Each stream is timestamped at the instant it is
read from its device, before any decoding or resizing, on a monotonic
clock. Stamping at read rather than at publication means the recorded time
reflects when the sensor was actually sampled, not when the processed
result happened to arrive.

\paragraph{Reference-time alignment.} When the recorder builds a frame it
fixes one reference time and samples every stream at that time. Images take
the nearest captured frame; proprioception and end-effector poses are
interpolated to the reference time, with orientations interpolated along
the shortest rotation. This replaces the naive rule of taking each stream's
latest value, which would place streams captured at different instants into
one frame as though they were simultaneous.

\paragraph{Drift budget.} For every stream and every frame the collector
records the residual between the reference time and that stream's nearest
real sample, and a live monitor shows these residuals during collection.
This matters because unsynchronised cameras cannot be re-phased in
software: their capture instants are fixed by their own clocks. The honest
deliverable is therefore a \emph{measured} drift, not a claim of perfect
synchronisation. We target a per-stream drift well within half a frame
period for the proprioceptive streams; the cross-camera residual is
whatever the hardware imposes, and the monitor makes it visible so an
operator can react to a starved bus. The same monitor also carries the
control prompts and a live progress readout --- the episodes recorded
against the session's target, together with the recorder's current phase, so
a save in flight is visible on the window --- so an operator drives a whole
session from the window rather than the terminal. \unjustified{} the half-frame target
is a working rule of thumb, not a value validated against downstream policy
performance.

\paragraph{Capture budget.} The bus, not only the clock, sets that residual.
Several colour streams share the host's controllers, and an uncompressed
transport does not fit: one uncompressed view at our resolution and rate needs
roughly eighteen megabytes per second, so a few of them exceed what the
controllers carry. The cameras then negotiate a fraction of the rate we ask
for. Nothing reports an error, because each frame that does arrive is valid;
the newest frame is simply already half an inter-arrival gap old, and the
measured residual reflects the capture rate rather than clock skew. We
therefore require a compressed transport on every stream, and a driver queue
short enough to bound staleness yet deep enough not to starve the camera. Both
halves are measured, not assumed. Requesting an uncompressed format held our
wrist cameras to under a third of their rate and left them a mean of
\SI{54}{\milli\second} behind the reference time; compressing them recovered
the full rate and cut that residual to \SI{20}{\milli\second}. A single-frame
queue looks attractive for freshness but halves the achieved rate, because the
driver has nowhere to place the next frame while the reader holds the only
buffer; two buffers sustain the full rate.

Fixing the transport does not exhaust the problem, because a camera's exposure
governs its rate as directly as the bus does: no camera can deliver frames
faster than it exposes them, so an exposure that grows past the frame period
forces the sensor to a slower one. Left automatic, exposure grows exactly when
the scene is dim, and the wrist views are the dim ones by construction --- they
look down at a close workspace that the arm itself shadows. The rate then
depends on the room rather than on the configuration, which is the property that
makes it hard to notice: the same rig measured in daylight and in the evening
gives different answers, and neither reports a fault. On our cameras the
relation is steep. Held to an exposure short enough to fit inside the frame
period, the wrist streams sustain their full rate; one step longer costs a fifth
of it; and the value automatic exposure selected under collection lighting gave
two thirds of the frame period to exposure alone, leaving the stream at little
over half its rate --- which was precisely the rate those streams had been
recording at. We therefore fix the exposure of the streams whose scene is
predictably dark. The cost is slight, because a fixed value chosen to fit the
frame period reaches the same image brightness that automatic exposure settles
on in good light; what is bought is a capture rate that is a property of the
configuration rather than of the hour of the day. The value is chosen per camera
and by eye, because two wrist cameras on one rig are not lit alike --- a single
figure that suits one over-exposes the other --- and because the judgement being
made is about a picture. The tool that sets them shows each stream's achieved
rate beside its image, so the one adjustment that costs rate is visible as it is
being made.

\paragraph{Recorded stream set.} The dataset records the colour views --- the
overhead, wrist and fingertip cameras --- the twelve-dimensional
proprioceptive state (five joints and a gripper per arm), the executed target
(Section~\ref{sec:contract}), and, when a depth-capable overhead camera is
fitted, its aligned depth image. That depth is sixteen-bit and is stored
losslessly outside the colour-video pipeline, because a depth image is not
three-channel colour and a colour-video encoder would corrupt it; the depth
scale and camera intrinsics are stored alongside. The full-rate side log
additionally holds the measured and target end-effector poses, the raw and
filtered operator input, and the gripper signals. Proprioception, command
and target share one unit convention and one channel layout, so a reader
never has to reconcile two coordinate conventions.

Each episode also carries an identifier taken from the wall clock at which it
began. Its position in the dataset cannot serve as its name, because removing
one episode renumbers every later one, so the same position denotes a different
recording afterwards, and any note an operator made against a number silently
comes to mean something else. The identifier is fixed when recording starts and
travels with the episode through renumbering, which lets a session log, a
review decision or a defect report keep pointing at the recording it was
written about.

An episode is also committed to storage the moment it is accepted, rather than
at the end of the session. This is not housekeeping: a dataset library that
accumulates episode records in memory and writes them in batches leaves the
session unreadable for as long as it runs, and loses every record still held in
memory if the session ends abnormally~--- and a recording whose record never
reached storage cannot be trained on, however complete its frames and video are
on the drive. A session lasting hours, on a rig where a camera can leave the bus
without warning, is precisely the case where that loss is likely and expensive.
Committing at episode granularity bounds any interruption to the recording in
flight, and has the further benefit that a session can be reviewed while it is
still being collected, so an operator discovers a framing or lighting mistake
during the sitting rather than after it.

\paragraph{Phase flag and gating.} Each frame carries a flag stating
whether teleoperation drove it. Homing moves and other non-teleoperation
motion are thereby marked and can be excluded when a policy is trained, so
the policy does not learn to freeze from frames where the recorded intent
fell back to the measured state (Section~\ref{sec:contract}). Gating then acts
at the episode level, and it is graded rather than absolute. A consumer bus
camera occasionally disappears for a second or two and returns by itself, and
discarding a long, otherwise good demonstration for such a blink costs an
operator far more than the missing second does. The recorder therefore holds a
stream loss: it stops adding frames, waits, and resumes the same episode once
every stream delivers again, abandoning the episode only if the stream stays
away. A session fault still discards outright.

We do not let that recovery hide a defect. A held episode really does contain a
gap, and because frame timestamps follow the frame index, the saved episode
would otherwise present one frame period across a hole that lasted seconds. The
recorder counts every hold and reports the total with the episode, so an
operator can inspect the episode and remove it if the motion jumps across the
join. Any end the operator did not ask for also sounds an audible cue, because
an operator attending to the arms cannot see the terminal, and would otherwise
continue demonstrating into a recording that had already stopped.
\unjustified{} the waiting period is set from the observed camera
reopen interval, not from any measurement of how large a gap a policy tolerates.

\paragraph{Session consistency.} A dataset is usually gathered over several
sittings, so its stream configuration must not drift between them. When a
session extends an existing dataset, the recorder reads back which streams
that dataset already holds --- the colour cameras, whether depth is present,
whether the end-effector channels were stored, and the frame rate --- and
follows them, rather than trusting flags typed afresh each time. The operator
therefore names the dataset and the task; the schema comes from the dataset
itself, and every episode in it shares one definition.

\section{Teleoperation-to-autonomy contract}
\label{sec:contract}

A demonstration dataset is only usable for autonomy if the label a policy
learns to emit is the same quantity, in the same frame, that the robot will
later execute from that policy. This section states that principle, then
its joint-space and end-effector instances, and finally the frozen
constants that tie the two ends together.

\paragraph{The principle.} The recorded action is the command the robot
actually followed, not the raw operator input. Between operator and motors
sit a rate limiter, the joint limits, a deliberately capped gripper
mapping, and, for the Cartesian interface, the reachable-workspace
projection. We record the output of that chain. At inference the policy's
output must pass back through the same chain: a joint-space prediction
through the same limiter and gripper cap, an end-effector prediction
through the same solver and projection. If inference skipped these steps,
the executed motion would leave the distribution the policy was trained on.

\paragraph{Target, not measured pose.} The recorded intent is a goal the
arm is moving towards, and it is kept strictly separate from the measured
state, which is where the arm currently is. When no fresh target exists ---
during a scripted homing move, for example --- the label falls back to the
measured state for that frame, and the phase flag of
Section~\ref{sec:collection} marks the frame so it can be excluded. This
keeps the label defined at every frame without teaching motion where none
was intended.

\paragraph{Joint-space policy.} The joint-space label is the commanded
joint target for both arms and the commanded gripper opening. An autonomous
joint-space controller emits the same quantity and sends it through the
same limiter and gripper cap that produced the labels.

\paragraph{End-effector policy.} The end-effector label is the constrained
Cartesian target for each arm, expressed in that arm's own base frame as a
position and an orientation. An autonomous end-effector controller emits
this target and replays it through the same inverse kinematics and the same
envelope projection. Because that target only exists when the Cartesian
solver runs, the end-effector features are recorded only under the
head-mounted interface, not under direct joint teleoperation.

\paragraph{Frozen action definition.} The constants that define the action
--- the units, the joint order, the gripper cap, the input scales, the
end-effector frame convention and the orientation representation --- are
stored with each dataset. Training and inference read the same definition,
so they cannot silently disagree. The end-effector state and target are
stored under neutral keys that a standard policy loader ignores, so a
default joint-space policy is unaffected by their presence, while an
end-effector experiment maps the stored target onto its action explicitly.

\paragraph{Ablating a sensor without disturbing the rest.} A platform that
records several viewpoints invites the question of how much each one
contributes, and that question is only answerable if the comparison changes
one thing. We therefore derive a \emph{view} of a collected dataset: a
dataset in its own right that declares a chosen subset of the cameras while
sharing the recorded frames themselves with the source. The episodes, the
frames and the labels are the ones that were collected; only the declared
sensor set differs. Because the learner reads the sensor set from the
declaration, an undeclared camera is neither decoded nor presented to the
policy, so the ablation costs less to train rather than more, and no
training code learns about the experiment. Deriving a view rather than
editing the recording also leaves the collected data untouched, which
matters because these datasets are expensive to gather and are shared
between experiments.

\paragraph{One instruction per task.} A language-conditioned policy reads
the instruction attached to each frame, so an instruction retyped partway
through a collection session splits one task into two. The view settles
this where it settles the sensor set: the spelling covering the most frames
is kept and the remainder are rewritten onto it. Weight of evidence decides,
not order of appearance, because the variant registered first is not
necessarily the intended one. Policies that ignore language are unaffected,
which is why the correction belongs to the derived dataset rather than to
the collection procedure.

\paragraph{Pretrained policies and fixed sensor slots.} A policy finetuned
from a large pretrained base does not derive its inputs from the dataset:
it arrives with a fixed set of sensor slots, each carrying what the base
learned about that viewpoint. Such a policy joins the study by mapping each
recorded camera onto the slot that means the same thing, rather than by
renaming the recording. A slot left without a camera is the ablation
itself, and the base handles it as absent evidence rather than as an error.
This keeps one dataset, and one contract, serving policies that derive
their interface from the data and policies that impose their own.

\paragraph{When there is no slot that means the same thing.} That mapping
presumes a correspondence, and a tactile platform breaks the presumption
twice. The bases available to us were pretrained on exterior and wrist
viewpoints, so no slot has ever been shown a deformable-contact image; and we
record more cameras than the slots number, so some camera must go unseen
whatever we do. We take both as platform decisions to be recorded rather than
hidden. A camera whose viewpoint the base does know keeps that slot by name,
so what the base learned about it is not discarded; a camera with no
counterpart is assigned the next free slot in a fixed order, which we state
plainly is an arbitrary placement rather than a correspondence; and a request
naming more cameras than there are slots is refused rather than silently
trained on fewer, because a run that quietly drops an input reports as a
result about the inputs it did not have. The assignment is recorded per run,
so it can be varied deliberately: which of the recorded views reaches a
pretrained slot is then an experimental factor of the study rather than a
property of the code.

\paragraph{Fitting more views than a policy accepts.} A stricter case admits
no assignment at all. Some recent policies concatenate their visual inputs
into a single frame of fixed size, so the number of views is not merely
bounded but exactly determined --- two, for the model we tried. Choosing two
of our cameras would discard every fingertip but one. Instead we compose: the
four fingertip streams are tiled into one image feature, frame for frame, and
that feature occupies a single view beside the exterior camera. The policy
receives every tactile signal at the cost of resolution rather than at the
cost of three of the four contacts, and the tiling is part of the dataset
presented to the learner rather than a step inside it, so it is visible,
reproducible and available to any policy on the same terms. \unjustified{}
that trading resolution for coverage helps is a design judgement here; we have
not measured a composed input against a chosen-subset one.

\paragraph{Training parity with simulation.} Because a collected dataset is
an ordinary dataset in the standard format, a policy trains on it through
exactly the pipeline the simulated study uses --- the same learner, the same
policy families --- with no real-specific training code. A small
plumbing check trains a policy on a collected dataset for a handful of steps
and reloads the resulting checkpoint, confirming the real recordings feed the
learner and that the checkpoint is ready for on-robot inference before any
long training run is committed. \unjustified{} the check validates wiring, not
learned skill.

\paragraph{Where the forward pass runs.} The contract fixes what a policy emits
and how that quantity reaches the motors; it does not fix which machine computes
it. The distinction matters in practice, because the computer that drives the
platform is chosen for its sensors and its buses rather than for its
accelerator, and the larger policy families do not sit comfortably on it. We
therefore let a separate machine serve the forward pass while the platform keeps
the cameras, the arms, the control rate and every safety decision. The split
costs little, because these policies plan a whole chunk of actions from a single
observation: the link carries one observation per chunk rather than one per
control step, and the following chunk is requested while the current one is
still being executed, so the transfer overlaps motion instead of delaying it.
The conversion from predicted action to motor command is unchanged and still
runs on the platform, which is what leaves the contract intact. Because the arms
hold torque throughout, a chunk that fails to arrive is treated as a fault
rather than as a delay: the controller holds the last commanded target, warns
the operator, and then releases the arms. \unjustified{} the tolerable holding
period follows from the control rate and the chunk length, not from any
measurement of how a policy's behaviour degrades across an interruption.

\paragraph{Where the training runs.} The same separation applies before
deployment. A collection lives on the platform, and the accelerators do not:
they are a shared cluster or a single remote workstation, reached over the
operator's existing credentials. We therefore treat a training destination as
configuration rather than as code --- how work is queued there, where its
scratch lies, and what its accelerator has been \emph{measured} to carry ---
and everything downstream of choosing a run is identical wherever it lands.
What this buys is not convenience but the ability to refuse early. A request
is resolved against the dataset and the destination before anything is
copied, so a camera the recording does not contain, a view count the policy
cannot accept, or a batch size that machine has been measured to exhaust is
reported in seconds rather than discovered hours into a reservation. The
destination's ceilings are recorded only where they were measured, and a
machine with none simply refuses nothing --- which we prefer to a plausible
number that would refuse runs that fit and admit runs that do not. The
workstation the operator is sitting at is treated as one of these destinations
rather than as a special case, since the same queueing discipline applies to a
single accelerator wherever it sits; what differs is only that its capacity
cannot be written down in advance, because the configuration is shared across
machines, and so it is asked of the machine at the moment a run is proposed.
That question earns its place: an accelerator too small does not make training
fail, it makes it fall back to the processor and continue, which for a run of
this length is indistinguishable from progress for as long as anyone is willing
to wait.

\paragraph{Following a run that is elsewhere.} A reservation of this length
is not a thing to submit and forget, and the two failures worth catching early
are opposite in character: a loss that is not falling, and a machine that has
quietly gone away. Neither is visible from a queue's own report, which says
only that a job exists. We therefore read progress from the record the training
itself leaves --- its log --- rather than adding instrumentation to the training
code, and we discover the runs on a machine by looking for them rather than
listing what was launched from any one place, so that a run started from a
terminal, by a colleague, or before the launcher existed is followed on the same
terms as one started from the console. Two properties of that log shape the
reading. Its step counter is formatted for a human and is therefore rounded
above a thousand, so the axis of a training curve must be reconstructed from the
logging interval rather than read off the line; and a progress display, where
the environment leaves one enabled, carries the exact position while a batch
scheduler suppresses it entirely. A reader that assumes either is present
silently mis-draws the curve on the machine where it is not. Staleness is judged
against the clock of the machine holding the log, never the reader's, since the
two need not agree; a run whose log has stopped growing relative to its own
observed cadence is reported as stalled rather than as running, which is the
state a lost machine leaves behind.

\paragraph{Surviving an interruption.} A run of this length outlives the
availability of the machine underneath it. We observed a shared workstation
restart beneath a reservation twice in as many days; the training log in each
case ends mid-line with no exception, and the only surviving evidence of the
cause is the machine's own record of when it last booted. What made this costly
was not the restart but the recovery: a partially complete run directory was
treated as debris and removed, so a reservation that had run for hours began
again from nothing. The contract is now that an interruption costs one
checkpoint interval rather than a run --- a completed run is reused, a partial
one is resumed from its last checkpoint, and only a directory holding no
checkpoint at all is discarded. This has a consequence for the reading of
progress described above. A resumed run appends to the log it already had, and
the training framework rebuilds its progress display over the steps that
\emph{remain}, so the second attempt's positions and the ordinal of its logged
lines both restart. Segmenting the log by the resumption itself recovers the
true axis; inferring it from the displayed totals does not, because it shifts
the first attempt's points as well whenever the target has not changed.

\paragraph{Which streams the policy uses.} A platform that records five
viewpoints and a proprioceptive vector invites the question of what each is
worth, and that question has two answers which should not be conflated. One
asks what a policy \emph{could} learn from a given subset, and is answered by
training a policy per subset and comparing them; it is the more familiar
experiment and it costs a reservation per cell. The other asks what a policy
that has already been trained on everything \emph{actually uses}, and is
answered from the trained weights alone in minutes. The two can disagree, and
where they do the disagreement is the result: a stream whose removal a
retrained policy compensates for is not the same as a stream the trained policy
ignores.

We answer the second by intervention wherever we can, and by inference only
where intervention cannot reach. Replacing one stream and re-planning is
behaviour rather than a claim about behaviour, so it is the measurement the
others are scored against; because there is no neutral substitute for an image,
the substitute used is reported alongside every result rather than treated as a
detail. Gradient attribution adds a decomposition whose parts sum to the change
in the quantity attributed, which is what makes several streams' contributions
comparable as shares of one whole rather than as unrelated magnitudes. Attention
weights, where the architecture exposes them, are reported beside these and
never in place of them. On our own checkpoint the attention mass over each
camera sits within a small fraction of what that camera's share of the tokens
would give it regardless of content: across two hundred frames the five cameras
span a factor of $1.1$ in attention where the behavioural measurement spans a
factor of forty, and on individual frames the two orderings agree only weakly
and disagree outright on a third of them. Attention here therefore sorts the
streams tolerably when averaged and understates the differences between them by
more than an order of magnitude, which is not the same as measuring importance;
the gradient decomposition, by contrast, reproduces the behavioural ranking
closely. We report the departure from what token count alone would predict,
never the raw mass. Finally, because an attribution averaged over an episode answers
almost nothing --- a contact sensor cannot contribute while the hand is still
travelling --- contributions are read against phases segmented from the
recorded gripper channels themselves, with thresholds drawn from each episode's
own distribution rather than fixed in advance.

\section{Replicability}
\label{sec:replicability}

Reproducing the platform needs its measured geometry, its camera
calibration, its rates, and its action constants. This section lists what
is measured and stored, and the procedure that confirms the rig is ready.

\paragraph{Geometry and provenance.} The rig geometry derives from one
parametric computer-aided-design description, and the same description
generates the digital twin. The documented dimensions and the physical
build therefore cannot drift apart. The validation object used for simple
pick-and-place checks is a printed cube of stated side length, sized so the
capped gripper can grasp it.

\paragraph{Camera calibration.} The overhead depth camera's intrinsic
parameters and its depth scale are read from the device and stored with
every dataset, so the depth images are metrically interpretable after the
fact. The exterior cameras are on fixed mounts, and their poses relative to
the rig are measured and recorded. A camera that shifts between collection
and deployment breaks a policy silently, so a fixed and recorded pose is
part of the platform, not an afterthought.

\paragraph{Rates and timing.} The dataset frame rate, the full-rate side
log rate, and the proprioception rate are stated for each session, together
with the measured drift budget of Section~\ref{sec:collection}. These
numbers let a reader judge the temporal quality of the data without
rerunning the rig.

\paragraph{Action constants.} The frozen action definition of
Section~\ref{sec:contract} --- units, joint order, gripper cap, input
scales, end-effector frame and orientation convention --- travels with the
dataset. It is exactly what an autonomous controller must reproduce, so
storing it beside the data removes a common source of silent train and
inference disagreement.

\paragraph{Readiness procedure.} A preflight check confirms, before
collection, that every arm is calibrated and bound to a stable alias, that
the wrist joints are limited to protect the fingertip cables, that the arms
are at their home pose, that every camera and the depth device open at
their configured resolution, that the dataset volume has headroom, and that
the host-tuning items which reduce timing jitter are set. It reports one
ready-or-not verdict, so an operator does not begin a long session on a rig
that is quietly misconfigured.

\section{Difficulties}
\label{sec:difficulties}

This section records the obstacles met while building the platform and how
each was resolved, so the result is reproducible and its remaining weak
points are visible. Each obstacle gets one short subsection.

\paragraph{Temporal skew across streams.} Taking each stream's latest value
at frame time placed samples captured at different instants into one frame.
The resolution is to stamp every stream on read, align all streams to one
reference time per frame, and log the residual drift
(Section~\ref{sec:collection}). A residual remains: unsynchronised cameras
keep an irreducible cross-camera skew that we measure but cannot remove in
software, because their capture clocks are independent.

\paragraph{The stale-command labelling trap.} When no fresh command exists,
the recorded action falls back to the measured state. On a moving homing
frame this would teach a spurious command to hold still. The resolution is
the per-frame teleoperation flag, which marks such frames for exclusion,
together with a logged count of how often the fallback fired, so the
problem is visible per episode rather than hidden in the data.

\paragraph{Depth cannot use the colour pipeline.} Sixteen-bit depth pushed
through a three-channel colour-video encoder is quantised and corrupted.
The resolution is to store depth losslessly as a per-frame image beside the
dataset, keyed to the dataset frame index, with its scale and intrinsics
recorded alongside.

\paragraph{Shared-bus bandwidth.} Several compressed image streams on
shared buses can starve one another and drop frames. The resolution is
compressed capture, a per-stream rate and dropped-frame counter shown live,
and episode-level discard when a stream is lost. \unjustified{} the exact
number of simultaneous streams a given host and hub can sustain is
machine-specific; we measure it per rig rather than derive it.

\paragraph{Appearance drift.} Automatic exposure and white balance make the
same scene look different across episodes and between collection and
deployment. The resolution is to lock exposure where the device allows and
to record the chosen settings, so appearance is a controlled and documented
variable rather than a hidden one.

\paragraph{Host timing jitter.} Background processes and power-saving
governors cause the recorder's frame loop to overrun its budget. The
resolution is the readiness check's host-tuning items and the recorder's
own overrun warnings, which surface a slow tick immediately. \unjustified{}
the overrun and staleness thresholds are set by convention and have not
been tuned against downstream policy performance.

\paragraph{Coverage for autonomy.} Clean expert trajectories
under-represent the off-nominal and recovery states an autonomous policy
inevitably visits, where small errors compound~\citep{ross2011dagger}. The
resolution is a collection protocol that deliberately varies object
placement and includes recoveries from near-failures, in the spirit of the
low-cost bimanual demonstration collection of \citet{zhao2023aloha}. This is a data-design obligation on the
operator, not something the recording software can supply.

\section{Conclusion and ongoing work}
\label{sec:conclusion}

We have described a real-world data-collection platform whose design is
governed by a single requirement: the recorded data must let a policy run
autonomously on the same rig. Two commitments follow from that requirement.
Every stream is timestamped on read and aligned to one reference time per
frame, with the residual drift logged, so temporal quality is measured
rather than assumed. And the recorded action is the constrained command the
robot actually executed, stored in both joint and end-effector form, with a
contract that inference replays a policy's output through the identical
command path.

The platform is measured and its defining constants travel with each
dataset, a readiness check confirms the rig is at a fixed collection state,
and the difficulties section keeps an honest record of what remains
imperfect --- chiefly the irreducible cross-camera skew and the several
thresholds still set by convention.

This is a living document. An experiments section will report the first
datasets collected on the platform, the measured drift budgets achieved,
and the autonomous performance of a joint-space and an end-effector policy
trained from the same episodes, once those runs exist.


\bibliographystyle{plainnat}
\bibliography{references}

\end{document}
